\documentclass[11pt]{article}
\usepackage[margin=1.1in]{geometry}
\usepackage[cmex10]{amsmath}
\usepackage{amssymb,amsthm}
\usepackage[T1]{fontenc}
\usepackage[hidelinks,hypertexnames=false]{hyperref}
\allowdisplaybreaks[2]
\setlength{\emergencystretch}{2em}
\title{A quadratic row bound for balanced arrays}
\author{ProgramComputer}
\date{}
\newtheorem{theorem}{Theorem}
\newtheorem{corollary}[theorem]{Corollary}
\DeclareFontShape{OMX}{cmex}{m}{n}{<-> cmex10}{}
\begin{document}
\maketitle

\begin{abstract}
We study the polynomial counting labeled binary $n\times d$ arrays with row sums
at most $r$ by their common column sum. The theorem gives one
absolute $C$ such that for every $d\ge2$, every $1\le r\le d$, and $n\ge Cd^2$,
all its zeros are simple and strictly negative. Together with the rank-dependent
dimension bound, this covers all dimensions for large $n$ when $r\le c\log n$
for a sufficiently small absolute $c>0$.
\end{abstract}

\section{Statement}
The \href{https://www.symmetricfunctions.com/realRootedOpenProblems.htm\#conjBalancedArrays}{balanced-array conjecture of Per Alexandersson}
asks for real-rootedness throughout the parameter range. Let $a_k$ count the
arrays in the abstract with each column sum equal to $k$, and set
\[
P_n^{(r,d)}(q)=\sum_{k=0}^{\lfloor nr/d\rfloor}a_kq^k.
\]
\begin{theorem}
There is an absolute $C>0$ such that for every integer $d\ge2$, every
$1\le r\le d$, and every integer $n\ge Cd^2$, the polynomial $P_n^{(r,d)}$
has exactly $\lfloor nr/d\rfloor$ simple, strictly negative zeros.
\end{theorem}
\begin{corollary}
There are absolute $c>0$ and $N$ such that for every $n\ge N$ and every $d\ge2$,
the same conclusion holds whenever $1\le r\le\min(d,c\log n)$.
\end{corollary}
The constants $C$, $c$ and $N$ are existential. At $d=1$, the polynomial is $(1+q)^n$, with a zero of multiplicity $n$ for $n\ge1$. The original all-parameter conjecture remains open.

For the original real-rootedness conjecture, any infinite sequence of
distinct counterexamples has $n\to\infty$: when $n$ is bounded,
\href{https://github.com/ProgramComputer/open-math-research/blob/04564a728c21eab6ce511b98a48ccaac754648ff/topics/balanced-arrays/orientation-normalization.md}{coefficient separation}
covers $d\ge n$, leaving finitely many triples. Along such a sequence the
cap is eventually greater than $c\log n$, and the dimension lies between
constant multiples of $\sqrt n$ and $n/\log n$. The upper restriction uses
the previous \href{https://github.com/ProgramComputer/open-math-research/tree/04564a728c21eab6ce511b98a48ccaac754648ff/topics/balanced-arrays/dense-regime}{high-dimensional bound}.

\section{Reduction to a sequence}

Write $P(q)=\sum_{k=0}^D a_kq^k$ and $T=nr/d$, $D=\lfloor T\rfloor$. For every integer $0\le k\le D$, fill $dk$ slots split into $n$ groups of at most $r$ with the repeated sequence of columns $1,\ldots,d$. This realizes a valid array because no row repeats a column and the total is divisible by $d$. Thus all coefficients on $[0,D]$ are positive.

If the theorem admitted no such $C$, we could select failures with $n/d^2\to\infty$. If $d$ were bounded along a subsequence, finitely many caps would be involved and the fixed-cap eventual theorem would contradict these failures. The unrestricted case $r=d\ge2$ is already covered by the binomial multiplier theorem. We may assume $d\to\infty$ and $1\le r<d$, allowing $r$ to vary. Set $b=d^2/n\to0$, $\phi=\pi/d$ and $\nu=(d-1)/2$. Then $D/d\to\infty$. \href{https://github.com/ProgramComputer/open-math-research/tree/f7e10f84dd21fec44ca58c60bae0879d12ebd89b/topics/balanced-arrays/fixed-cap}{Fixed-cap theorem} \enspace \href{https://github.com/ProgramComputer/open-math-research/blob/04564a728c21eab6ce511b98a48ccaac754648ff/topics/balanced-arrays/finite-linear-reduction.md}{Unrestricted case}

\section{Choosing and locating the phase levels}

Under a positive row tilt $w$, the size $S$ has weights $\binom ds w^s$ for $0\le s\le r$. Write $H(w)$ for their sum, $\mu=\mathbb ES$, $V=\mathrm{Var}(S)$, $h=r-\mu$, $y=n\mu/d$ and $G(s)=s(d-s)/(d(d-1))$. Let $w_x$ denote the tilt with $y=x$, unique on $0<x<T$ because the derivative of the mean in $\log w$ equals $V>0$. Write

\begin{equation*}
\lambda_t=\frac{\mathbb E[G(S)e^{itS}]}{\mathbb E e^{itS}},
\qquad \Psi(w)=n\arg H(we^{i\phi})-\nu\arg\lambda_\phi .
\end{equation*}

Take arguments continuously from angle zero. Both positive sums defining the arguments have angles in $[0,r\phi]$ with $r\phi<\pi$, so they are nonzero. As $w$ increases from zero to infinity, $\Psi$ has respective limits $-\nu\phi$ and $T\pi$, using the first nonzero and last degrees of the numerator defining $\lambda_\phi$. Every internal level $k\pi$ therefore has a first hitting point $v_k$, and continuity gives $v_k<v_{k+1}$.

We use two estimates from the derivation of the quartic-log bound that hold without its height assumption. For every tilt,

\begin{equation}
|\Psi(w)/\pi-y|\le C_0 d+C_0\min(y,T-y)/\sqrt d .
\tag{M1}
\end{equation}

Indeed the variance bounds $V\le d/4$, $V\le\mu(1+\mu)$ and $V\le h(1+h)$ imply $V\le C\sqrt d\min(\mu,h)$. Taylor's inequality gives $|\arg H(we^{i\phi})-\mu\phi|\le C V/d^2$, and $|\arg\lambda_\phi|\le\pi$. The cited variance bounds follow respectively from the success/failure size-bias comparison for the truncated binomial and from log-concavity of the tail probabilities of $S$ and $r-S$. \href{https://github.com/ProgramComputer/open-math-research/blob/04564a728c21eab6ce511b98a48ccaac754648ff/topics/balanced-arrays/very-tall/proof.md}{Phase estimates and proof of the variance bounds}

At $k/d\to c_0\in(0,\infty)$, evaluated at $v_k$, (M1) first gives $y=O(d)$, hence $\mu=O(b)$. Putting $u=dw$ gives $u\le\mu/(1-\mu)$ and size weights bounded above by $u^s/s!$. Therefore $\mathbb E|S-\mu|^3=O(\mu)$, and reweighting with $G$ gives expected distance $O(\mu)$ from size one. The latter follows from $G(s)/G(1)\le s$ and $\mathbb P(S=1)\asymp\mu$. Taking the imaginary part after centering at the mean in the characteristic function gives a radial angular error $O(\mu\phi^3)$, and $\arg\lambda_\phi=\phi+O(\mu\phi)$. Since $n\mu\phi^3=O(y/d^2)=o(1)$, the level equation yields $y(v_k)=k+1/2+o(1)$. On bounded translates of $k$, $V_x\sim dx/n$ uniformly and $\frac{\mathrm d}{\mathrm dx}(d\log w_x)=b/V_x$. Integration thus shows

\begin{equation}
d\log v_k=d\log w_k+d/(2k)+o(1).
\tag{M2}
\end{equation}

If $k/d\to\infty$, (M1) ensures $y(v_k)/d\to\infty$. This implication requires no assumption about the distance from the upper endpoint.

\section{Signs and order}

Put $t_k=\sqrt{a_{k-1}/a_{k+1}}$ and $\Gamma_k=a_k^2/(a_{k-1}a_{k+1})$. Use the radii $t_k$ for $1\le k\le d$ and $v_k^d$ for $d<k<D$. There is only one change of radii. The upper angular estimate treats every internal level near saturation, including the levels within a bounded distance of the degree.

For $k/d\to0$, the orientation inequality gives $\Gamma_k\ge(1+1/k)^d\to\infty$ and strict log-concavity throughout the support. At $t_k$, normalize by $(-1)^ka_kt_k^k$. The two neighbors have magnitude $1/\sqrt{\Gamma_k}$, and the remaining alternating tails start with a positive term and have decreasing magnitudes. This gives a strictly positive limit inferior. \href{https://github.com/ProgramComputer/open-math-research/blob/04564a728c21eab6ce511b98a48ccaac754648ff/topics/balanced-arrays/orientation-normalization.md}{Orientation inequality}

At $k/d\to c_0\in(0,\infty)$, the lower-quota comparison proves $\log\Gamma_{k+j}\to1/c_0$ at every fixed integer translate, and $\log t_k=d\log w_k+d/(2k)+o(1)$. This uses exact finite differences of the margin normalizing factors and controls the sparse margin correction after conditioning. At $x=t_k$ or $x=v_k^d$ we therefore have $\log(x/t_k)=o(1)$ by (M2). For fixed $j$, $a_{k+j}x^j/a_k\to e^{-j^2/(2c_0)}$ by multiplication of adjacent ratios starting with the two neighbors. Both distances to the support endpoints tend to infinity. By log-concavity the outward consecutive ratios after any large fixed distance $L$ are bounded by a geometric factor less than one, using their limits at $L$. Sending $L$ to infinity gives

\begin{equation*}
\frac{(-1)^kP(-x)}{a_k x^k}\longrightarrow
\sum_{j\in\mathbb Z}(-1)^j e^{-j^2/(2c_0)}
=\sqrt{2\pi c_0}\sum_{m\in\mathbb Z}e^{-2\pi^2 c_0(m+1/2)^2}>0
\end{equation*}

by Poisson summation. \hyperref[sec:quadratic-lower-quota]{Lower-quota comparison}

Now let $k/d\to\infty$ and set $w=v_k$. Integrating the row polynomial power on the torus $|z_i|=1$, $\prod_i z_i=-1$ gives $P(-w^d)$, since only common column exponents survive. Choose the fixed constant $h_0>0$ from the near-full angular lemma. If $\mu\le1/10$, then $n\mu/d^2=y/d\to\infty$ and the small-mean integral estimate applies. Its only size condition is $n\mu/d^2\to\infty$. If $\mu\ge1/10$ and $h\ge h_0$, the quadratic interior estimate applies. Each gives two conjugate principal terms with sum equal to twice their positive amplitude times $\cos\Psi(w)$, with error $o(1)$ on the scale of one term. Thus the sign is $(-1)^k$. \href{https://github.com/ProgramComputer/open-math-research/blob/04564a728c21eab6ce511b98a48ccaac754648ff/topics/balanced-arrays/very-tall/very-tall-nearly-empty.md}{Small-mean estimate, first section} \enspace \hyperref[sec:quadratic-interior-saddle]{Quadratic interior estimate}

For $h\le h_0$ the near-full angular lemma gives the same sign at every internal phase level. In this range the equality $\Psi=k\pi$ yields $T-k=(nh/d)(1+O(d^{-2}+d/n))$, so $nh$ is at least of order $d$ even at the last internal level. If $nh/d^2\to\infty$, just the two principal terms determine the sign. When $nh/d^2$ is bounded, an odd multiple $\ell\phi$ with $|\ell\phi|\le M/\sqrt{nh}$ contributes a phase equal to $\ell k\pi+o(1)$ for fixed $M$. All such terms have the required sign, and their summed positive amplitudes dominate the remaining Gaussian tail for a sufficiently large $M$. This covers both integral and nonintegral $T$. \hyperref[sec:near-full-all-angular-signs]{Near-full sign theorem}

The geometric points are increasing by strict log-concavity. The phase points are increasing by their first-hitting definition. At the join, $\log(t_{d+1}/t_d)=(\log\Gamma_d+\log\Gamma_{d+1})/2\to1$, and at index $d+1$ the replacement by the phase point changes its logarithm by $o(1)$. So all test points are strictly ordered for large $d$. Use zero and positive infinity for the endpoint signs $1$ and $(-1)^D$. If an interior sign failed along a sequence, a subsequence would have $k/d$ tending to zero, a positive finite number or infinity. The three cases exclude such a failure. The alternating signs give $D$ distinct negative zeros, proving the theorem.

\section{Ranks up to a constant times the logarithm}

Fix $\varepsilon=1/10$ in the rank-dependent column bound and let $B$ denote its dimension constant. If $n\ge N$ and $r\le c\log n$, choose $N$ sufficiently large that $r\le n^{1/10}$. If $n\ge Cd^2$, the quadratic row theorem applies. Otherwise $d^2>n/C$, which implies $d\ge B\sqrt{nr/\log n}$ whenever $c\le 1/(CB^2)$. The rank-dependent bound therefore applies in the other range. The general conjecture remains open. \href{https://github.com/ProgramComputer/open-math-research/tree/ba7ddfd7efb37cac33769f4fdfed749e1b5bf58a/topics/balanced-arrays/dense-regime/rank-dependent}{Rank-dependent theorem}


\appendix
\clearpage
\section{Lower quota coefficients at the quadratic row scale}\label{sec:quadratic-lower-quota}

Let $d\to\infty$, $n/d^2\to\infty$, $1\le r<d$ and $k/d\to c\in(0,\infty)$. With the original labeled balanced-array counts $a_i$, set $\Gamma_i=a_i^2/(a_{i-1}a_{i+1})$ and $t_i=\sqrt{a_{i-1}/a_{i+1}}$. Let $H(w)=\sum_{s=0}^r\binom ds w^s$ and let $w_x$ be the positive tilt whose size mean is $dx/n$. Then, for every fixed integer $j$,

\begin{equation*}
\log\Gamma_{k+j}\to 1/c,\qquad \log t_k=d\log w_k+d/(2k)+o(1).
\end{equation*}

The cap may grow arbitrarily with the dimension. The following proof separates caps at $L=\lceil\log d\rceil$.

\subsection{Caps at or above the cutoff}

At every $i$ in a fixed window about $k$, put $M=di$, $p=i/n$ and $\mu=M/n$. Choose each column independently as a uniform $i$-subset of the rows. The probability of a specified row containing a specified $(L+1)$-subset of columns is $p^{L+1}$. The expected number of such subsets in the array is at most $n\mu^{L+1}/(L+1)!=M\mu^L/(L+1)!=o(1)$, since $M=O(d^2)$, $\mu\to0$ and $L\ge\log d$. Every array violating the cap has such a subset. Hence $a_i=\binom ni^d(1+o(1))$ throughout this window.

To compare the tilt, put $u=dw_k$. Markov's inequality and the ratio of size-one and size-zero masses give $u\le\mu_k/(1-\mu_k)$, where $\mu_k=dk/n$. Under the uncapped law at $w_k$, the mean $m=dw_k/(1+w_k)$ exceeds $\mu_k$ by at most $\sum_{s>r}s u^s/s!\le 2u^{L+1}/L!$ for large $d$. Indeed truncation deletes the upper tail, and $m-\mu_k$ is the unconditional tail expectation of $S-\mu_k$. Thus $d(m-\mu_k)/\mu_k=o(1)$. Because the inverse uncapped tilt is $m/(d-m)$, this gives $d\log(w_k/(k/(n-k)))=o(1)$. The adjacent ratios of the unrestricted binomials now give both conclusions. Explicitly their centered log ratio is $d\log(k/(n-k))+(d/2)\log(1+1/k)-(d/2)\log(1+1/(n-k))$, and their log curvature at $i$ is $d\log(1+1/i)+d\log(1+1/(n-i))$.

\subsection{Caps below the cutoff}

For an integer $i=k+O(1)$ let $X$ have tilted sizes with weights proportional to $\binom ds w_i^s$. Uniformly, $dw_i=\mu+O(\mu^2)$, its variance $V=\mu+O(\mu^2)$, $\mathbb E[X(X-1)]=O(\mu^2)$ and $\mathbb E|X-\mu|^3=O(\mu)$. These follow by bounding size weights by $(dw_i)^s/s!$ with the first two retained. The masses at zero and one have product comparable to $\mu$. Hence
$|\mathbb E e^{itX}|\le e^{-c_0\mu t^2}$ for $|t|\le\pi$. Fourier inversion at the integer mean gives $\mathbb P(T=M)\sim(2\pi nV)^{-1/2}$ for $T=\sum_{j=1}^n X_j$, by independence, the third-moment Taylor approximation at $t=O((nV)^{-1/2})$, and Gaussian domination. Here $M\asymp d^2\to\infty$.

For a feasible row profile $s_j\le r$ summing to $M$, put $U=\sum_j s_j(s_j-1)$ and $R=\sum_j(s_j-\mu)^2=M(1-\mu)+U$. Greenhill--McKay--Wang's sparse binary margin formula gives its count as

\begin{equation*}
K_i\prod_j\binom d{s_j}\,
 \exp\!\left(-\frac12+\frac{R}{2M(1-p)}+o(1)\right),
\qquad K_i=\frac{\binom ni^d}{\binom{nd}{di}},
\end{equation*}

uniformly over profiles. Indeed the maximum degree product is $ri=O(d\log d)=o(M^{2/3})$, normalized column variance is zero, and the normalized row variance $R/(M(1-p))\le r/(1-p)=O(\log d)=O(M^{1/3})$. These are the hypotheses of Corollary 5.2 under Theorem 1.3. Each such profile is feasible by cyclically assigning the $M$ column labels to consecutive row slots. \href{https://web.maths.unsw.edu.au/~csg/papers/irregular-final.pdf#page=33}{Authors' final text, Corollary 5.2} \enspace \href{https://doi.org/10.1016/j.jcta.2005.03.005}{Journal DOI}

The law of the profile given $T=M$ weights it in proportion to $\prod\binom d{s_j}$. Write $W(s)=\exp(s(s-1)/(2M(1-p)))$ and $Z=\mathbb E W(X)$. Since $s^2/M=O((\log d)^2/d^2)=o(1)$, we have $1\le Z\le1+C\mu^2/M$. Thus $Z^n=1+O(\mu)$. We need the same upper bound after conditioning, since the crude maximum of $U/M$ could grow. The law reweighted by $W/Z$ still has product of zero and one masses comparable to $\mu$, as $W(0)=W(1)=1$. For any mixture of $n-1$ independent rows, each with either the original law or this reweighted law, the supremum of the total-mass probabilities is $O(M^{-1/2})$ by the same Fourier envelope. Telescope the product minus one, marking row $j$, weighting all previous rows by $W$ and the marked row by $W-1$. Conditional on the marked row, normalize the previous $j-1$ rows and apply the mixed supremum bound to all remaining $n-1$ rows. This gives

\begin{equation*}
0\le\mathbb E[(\prod_{j=1}^n W(X_j)-1)\mathbf1_{T=M}]
\le C M^{-1/2}\sum_{j=1}^n Z^{j-1}(Z-1)
=O(M^{-1/2}\mu).
\end{equation*}

Dividing by $\mathbb P(T=M)\asymp M^{-1/2}$ proves the conditional claim. The deterministic part of the margin exponent is $(p-\mu)/(2(1-p))=o(1)$. Therefore $a_i=K_i A_i(1+o(1))$ with $A_i=[w^{di}]H(w)^n$.

For completeness, $\Delta_i=\log K_i-\log K_{i-1}$ is exactly

\begin{equation*}
\sum_{j=1}^d\left(\log\frac{n-i+1}{n-i+j/d}
                        +\log\frac{i-1+j/d}{i}\right)
=-\frac{d-1}{2i}+O(d/n+d/i^2).
\end{equation*}

This follows by expanding the logarithms and summing $1-j/d$. The $K_i$ term thus contributes $d/(2k)+o(1)$ to the signed centered first difference and $o(1)$ to each centered second difference. At real quota $x$ put $E(x)=n\log H(w_x)-dx\log w_x$. Then $E'(x)=-d\log w_x$ and $E''(x)=-d^2/(n V_x)\to-1/c$ uniformly on bounded translates of $k$. Its third derivative is $O(1/d)$ there, since $|dV/d\mu|\le \mathbb E|X-\mu|^3/V=O(1)$. The exact tilt formula says $A_i=e^{E(i)}\mathbb P(T=M)$, and the neighboring central probabilities have ratios tending to one by the local limit and $V_i\sim di/n$. Applying centered first and second differences gives the asserted formulas.

\section{Uniform interior saddle at the quadratic row scale}\label{sec:quadratic-interior-saddle}

Let $d\to\infty$, $1\le r<d$, and $v>0$. Let $S$ have probabilities proportional to $\binom ds v^s$ on $0\le s\le r$, and given $S$, let $I$ be a uniformly random $S$-subset of $[d]$. Put $H=\sum_{s\le r}\binom ds v^s$, $\mu=\mathbb ES$, $h=r-\mu$, $V=\mathrm{Var}(S)$, $\phi(t)=\mathbb E e^{itS}$, and $R(\theta)=\mathbb E e^{i\sum_{j\in I}\theta_j}$. Define $a_s=s(d-s)/[d(d-1)]$, $a=\mathbb E a_S$, $a(t)=\mathbb E[a_S e^{itS}]/\phi(t)$ and $t_*=\pi/d$. Fix $h_0>0$. Uniformly for $\mu\ge1/10$, $h\ge h_0$, and $n/d^2\to\infty$, one has

\begin{equation*}
\frac{P_n^{(r,d)}(-v^d)}{H^n}
=\frac{2}{\sqrt d}(2\pi n)^{-(d-1)/2}\,
\mathrm{Re}\!\left(\phi(t_*)^n a(t_*)^{-(d-1)/2}\right)
+o\!\left(\frac{|\phi(t_*)|^n}{\sqrt d}(2\pi n|a(t_*)|)^{-(d-1)/2}\right).
\end{equation*}

The error is measured on the size of a single saddle. Constants may depend on $h_0$. This is an interior statement; no endpoint or root count follows directly.

\subsection{Uniform moments and diagonal loss}

Use the success/failure bias comparison and modal-atom proofs in \href{https://github.com/ProgramComputer/open-math-research/blob/04564a728c21eab6ce511b98a48ccaac754648ff/topics/balanced-arrays/very-tall/very-tall-interior-saddle.md}{the fixed-cutoff saddle note}. They give $c_{h_0}\le V\le\mu(1-\mu/d)\le d/4$ and $a\ge \max(V/d,1/(11d))$.

The symmetric law of $\Delta=S-S'$ for an independent copy $S'$ is log-concave, by convolution. We make the tail bound uniform in $V$. Its one-sided survival $T(j)=\mathbb P(\Delta\ge j)$ for integers $j\ge0$ is log-concave: $T(j+1)/T(j)$ is the weighted average of successive mass ratios at indices at least $j$, and those ratios decrease in the index. Omitting the first index decreases this average. Thus $T(j+m)/T(j)\le T(m)/T(0)$. With $m=\lceil 4\sqrt{1+V}\rceil$, Chebyshev and $T(0)\ge1/2$ give $T(m)/T(0)\le1/4$. Iterating and using symmetry yields $\mathbb P(|\Delta|>x)\le C e^{-c x/\sqrt{1+V}}$. In particular $\mathbb E\Delta^4\le C(1+V)^2\le C_{h_0}V(1+V)$. Jensen conditional on $S$ gives $\mathbb E|S-\mu|^j\le\mathbb E|\Delta|^j$ for all $j\ge1$. The tail contribution to $\mathbb E\Delta^2$ beyond $L\sqrt{1+V}$ can be made less than any specified fraction of $V$ by taking $L$ large in terms of $h_0$.

As in the \href{https://github.com/ProgramComputer/open-math-research/blob/04564a728c21eab6ce511b98a48ccaac754648ff/topics/balanced-arrays/very-tall/very-tall-interior-saddle.md}{fixed-cutoff note}, decompose the symmetric decreasing law of $\Delta$ into uniforms on integer intervals $[-j,j]$. Uniforms with $j>2L\sqrt{1+V}$ contribute only a small fraction of the second moment, since a constant fraction of their second moment lies past $L\sqrt{1+V}$. The cosine loss of each uniform is at least $c\min(j^2t^2,1)$ for $|t|\le\pi$. Thus

\begin{equation*}
1-|\phi(t)|^2\ge c_{h_0}V\min(t^2,(1+V)^{-1}).
\end{equation*}

For $0\le t\le t_0=c_0/\sqrt{1+V}$, differentiating the cosine loss gives $\mathbb E[\Delta\sin(t\Delta)]\ge 2Vt-\mathbb E[\Delta^4]t^3/6\ge Vt$ for a sufficiently small $c_0$ depending on $h_0$. Also $1-|\phi(t)|^2\le Vt^2\le c_0^2$. In particular for $t_*\le |t|\le t_0$,

\begin{equation*}
|\phi(t)|/|\phi(t_*)|\le\exp[-c V(t^2-t_*^2)].
\end{equation*}

At $|t|\ge t_0$ there is a constant modulus gap from one, depending on $h_0$.

\subsection{Charts and expansion}

The negative evaluation is the Haar integral of $R(\theta)^n$ on the translated torus $\sum_j\theta_j=\pi\bmod 2\pi$. Its Euclidean covolume is $\sqrt d(2\pi)^{d-1}$ by the basis $2\pi(e_j-e_d)$ of its period lattice. Use disjoint charts at the $d$ diagonal angles $t=(2j+1)\pi/d\bmod 2\pi$, with $y\in \mathbf1^\perp$ and $\|y\|\le\eta$ for a small fixed $\eta>0$. Disjointness follows because two distinct diagonal points differ modulo $2\pi$ by a phase vector whose zero-sum lift has at least two coordinates differing by $2\pi$. An overlap would give a lift of norm at most $2\eta$. Using the nearest-anchor chord argument from the \href{https://github.com/ProgramComputer/open-math-research/blob/04564a728c21eab6ce511b98a48ccaac754648ff/topics/balanced-arrays/very-tall/very-tall-interior-saddle.md}{fixed-cutoff note}, the complement of the charts has $D_z=\sum_j|e^{i\theta_j}-d^{-1}\sum_i e^{i\theta_i}|^2\ge c\eta^2$.

The disjoint-pair bound of that note gives a fixed-size modulus loss at least $c\min(a_sD_z,1)$, since $\min(s,d-s)/d$ is comparable to $a_s$ for every size. The $O(1/d)$ Lipschitz bound of $a_s$ gives $\mathbb E a_S^2\le a^2+C V/d^2\le C'a^2$ by $a\ge \max(V/d,1/(11d))$. Cauchy--Schwarz yields $\mathbb P(a_S\ge a/2)\ge1/(4C')$. Averaging the fixed-size modulus losses therefore gives
$|R|\le \exp[-c\min(aD_z,1)]$.

Where $aD_z>1$, the bound $e^{-cn}$ is negligible relative to the principal size, since $n\gg d\log n$ and $\log(1/|\phi(t_*)|)=O(V/d^2)=o(1)$. On the remainder of the chart complement split the bound $e^{-cnaD_z}$ into two factors, using $D_z\ge c\eta^2$ in one. For the other partition by a closest anchor $j$ to the mean phase. Since $\sum_i|z_i-z_j|^2\le2D_z$, in principal relative angles $u_i=\arg(z_i/z_j)$ it is bounded by $e^{-c'na\sum_{i\ne j}u_i^2}$. The pushforward of Haar on the translated torus by relative angles at a fixed anchor is normalized Haar on the full $(d-1)$-torus. Indeed this map is the restriction of a continuous group homomorphism up to translation, and is surjective because $d\theta_j=\pi-\sum_{i\ne j}u_i\bmod2\pi$ has a solution for any angles. Integrating the bound for each of the $d$ anchors shows that this part is at most
$d(C/\sqrt{na})^{d-1}e^{-c_\eta na}$.
Dividing by the principal Gaussian size costs $e^{O(d)}$ and the principal radial modulus costs $e^{O(nV/d^2)}\le e^{O(na/d)}$. The relative error thus tends to zero when $n/d^2\to\infty$, because $na\ge n/(11d)\gg d$.

Write $Q_j=\sum_i y_i^j$ and $L=\sum_{j\in I}y_j$. Conditional on $S=s$, the first moment vanishes, the second is $a_sQ_2$, and the third is $\gamma_s Q_3$, where $\gamma_s=a_s(d-2s)/(d-2)$. The fourth-moment identity in the \href{https://github.com/ProgramComputer/open-math-research/blob/04564a728c21eab6ce511b98a48ccaac754648ff/topics/balanced-arrays/very-tall/very-tall-interior-saddle.md}{fixed-cutoff note}, applied to $m=\min(s,d-s)$ and using inclusion $p_1\le Ca_s$, $p_2\le C a_s^2$, gives $\mathbb E[L^4\mid S=s]\le C(a_sQ_4+a_s^2 Q_2^2)$. The functions $a_s,\gamma_s$ are $O(1/d)$-Lipschitz on $[0,d]$ for large $d$, and $|\gamma_s|\le Ca_s$. Therefore $\mathbb E a_S^2\le a^2+C V/d^2\le C a^2$, by combining $a\ge c/d$ and $a\ge V/d$. Centering $Z=S-\mu$, the same Lipschitz estimate gives

\begin{equation*}
\mathbb E[a_S|Z|]\le Ca\sqrt{1+V},\qquad
\mathbb E[a_S Z^2]\le C a(1+V).
\end{equation*}

Indeed bound $a_S\le a_\mu+C|Z|/d$, use $a_\mu=a+V/[d(d-1)]$, and use $\mathbb E|Z|^3\le(\mathbb EZ^2\,\mathbb EZ^4)^{1/2}\le C V\sqrt{1+V}$. Terms containing $V/d$ are bounded by $a$.

Throughout each chart, $|R(t\mathbf1+y)-\phi(t)|\le aQ_2/2$, by centering the linear term. For $|t|\ge t_0$, choosing $\eta$ small enough makes this perturbation less than half the uniform diagonal gap. Their normalized integral is negligible by $n\gg d\log n$.

Inside a near chart $|t|\le t_0$, set $\gamma(t)=\mathbb E[\gamma_S e^{itS}]/\phi(t)$. The preceding weighted moments yield $|\gamma(t)|\le Ca$, $|\mathrm{Im}\,\gamma(t)|\le Ca\sqrt{1+V}|t|$, and

\begin{equation*}
|\mathrm{Re}\,a(t)-a|\le Ca(1+V)t^2,\qquad |\mathrm{Im}\,a(t)|\le Ca\sqrt{1+V}|t|.
\end{equation*}

For the real bound, center the numerator and denominator at $\mu$, use the cosine quadratic bound for their real perturbations, and use the sine linear bounds for imaginary parts when taking the quotient. The denominator modulus is uniformly close to one for small fixed $c_0$. Choose $c_0$ small enough that $\alpha_t=\mathrm{Re}\,a(t)\asymp a$. Then Taylor and expansion of the logarithm on $\|y\|\le\eta$ give

\begin{equation*}
\log(R(t\mathbf1+y)/\phi(t))
=-\tfrac12a(t)Q_2-\tfrac i6\gamma(t)Q_3+O(aQ_4+a^2Q_2^2).
\end{equation*}

The log-square term is $O(a^2 Q_2^2)$ because $|R-\phi|\le aQ_2/2$.

\subsection{Uniform chart integrals}

Write $I_0=(2\pi/(na))^{(d-1)/2}$ for the unnormalized Gaussian integral. Taking moduli in a near chart leaves a Gaussian of coefficient $n\alpha_t$ and a positive exponential whose exponent is bounded by

\begin{equation*}
Cna\left(\sqrt{1+V}|t|\sum_j|y_j|^3+Q_4\right)+Cna^2Q_2^2.
\end{equation*}

By Cauchy--Schwarz it suffices to bound the expectations of the separated exponentials doubled. For the coordinate terms use $y_i=Z_i-\bar Z$ with independent $N(0,\sigma^2)$, $\sigma^2=(n\alpha_t)^{-1}$. Restrict additionally to $|\bar Z|\le\eta$, an independent event of probability tending to one. On the chart this forces $|Z_i|\le2\eta$ and $\sum|y_i|^j\le2^j\sum|Z_i|^j$. For each one-dimensional integral, $\exp(C'\sigma^{-2}[\sqrt{1+V}|t||Z_i|^3+Z_i^4])$ on this interval has expectation at most $1+C''(\sqrt{1+V}|t|\sigma+\sigma^2)$. This follows from $e^u-1\le u e^u$ and the doubled-variance Gaussian after choosing $\eta$ small, using $\sqrt{1+V}|t|\le c_0$. Multiplying gives a constant times $\exp(Cd[\sqrt{1+V}|t|\sigma+\sigma^2])$, at most $\exp(O(d(1+V)t^2+d/(na)+1))$ by Young's inequality.

For the radial term, $Q_2/\sigma^2$ has the $\chi^2_{d-1}$ law. On $Q_2\le Ld/(na)$ the doubled radial exponent is $O_L(d^2/n)=o(1)$. Above that threshold but with $\|y\|\le\eta$, absorb the doubled radial exponent using $Ca\eta^2$ as a small relative weakening of the Gaussian. The Gaussian normalization cost is at most $e^{C_1d}$ and the weakened tail is at most $e^{-cLd}$ for large fixed $L$ by the $\chi^2$ exponential moment. This contribution tends to zero for sufficiently large $L$.

Finally the ratio between Gaussian integral at $\alpha_t$ and $I_0$ is at most $\exp(Cd(1+V)t^2)$. Thus each near chart integral in modulus is at most

\begin{equation*}
C I_0|\phi(t)|^n\exp(Cd(1+V)t^2).
\end{equation*}

The principal Gaussian factor in modulus is comparable to $I_0$, since $|\mathrm{Re}\,a(t_*)/a-1|=O((1+V)/d^2)$ and $\mathrm{Im}\,a(t_*)/a=O(\sqrt{1+V}/d)$. For $|t|=j\pi/d$ at odd $j\ge3$, combine the last bound with the radial comparison. Since $nV\gg d(1+V)$, the relative bound is $\exp(-c n V(j^2-1)/d^2)$ up to an absolute constant. Summing over these phases gives $o(1)$.

At the principal chart the same exponential expectation bounds are uniform constants. For a real Gaussian of coefficient $n\alpha_{t_*}$ on the hyperplane, Wick and the zero sum give $\|Q_3\|_2=O(\sqrt d\,\sigma^3)$, as in the \href{https://github.com/ProgramComputer/open-math-research/blob/f7e10f84dd21fec44ca58c60bae0879d12ebd89b/topics/balanced-arrays/fixed-cap/tall-sparse-transverse-saddle.md}{fixed-cap saddle calculations}. The triangle inequality for moments gives $\|Q_4\|_2=O(d\sigma^4)$ and $\|Q_2^2\|_2=O(d^2\sigma^4)$. Thus the $L^2$ size of the log correction to the full complex quadratic, multiplied by $n$, is $O(d/\sqrt n+d^2/n)$, using $a\ge1/(11d)$. Apply Cauchy--Schwarz to $|e^B-1|\le|B|e^{(\mathrm{Re} B)_+}$. The Gaussian tail beyond $\|y\|=\eta$ is negligible since $na\gg d$. The real Gaussian normalization and absolute complex Gaussian normalization are comparable by the principal estimates above. This proves the principal local relative $o(1)$ and the two-saddle formula after dividing by the torus covolume.


\subsection{Quantitative consequence for a large fixed quadratic constant}

With $h_0$ fixed, the proof also works for sign testing uniformly at $n\ge C_*d^2$ with $C_*$ sufficiently large depending on $h_0$. Relative to one saddle the local error is $O_{h_0}(C_*^{-1/2}+C_*^{-1})$ for large $d$, and the total near nonprincipal error is $O_{h_0}(e^{-c_{h_0}C_*})$, by the summable bound in odd $j$. In the radial $\chi^2$ estimate the central region contributes $e^{O_{h_0,L}(1/C_*)}$, still a bounded factor. Both far errors tend to zero as $d\to\infty$ for sufficiently large fixed $C_*$, since $na\ge C_*d/11$. Thus the absolute error on one-saddle scale can be made as small as desired by increasing $C_*$. An asymptotic relative $o(1)$ is asserted only when $n/d^2\to\infty$.

\section{All internal angular signs near the upper edge in growing dimension}\label{sec:near-full-all-angular-signs}

Let $d\to\infty$, $1\le r<d$, and $n/d^2\to\infty$. Put $H(w)=\sum_{s=0}^r\binom ds w^s$. For $t>0$, give the row size $S$ probabilities proportional to $\binom ds t^s$, and conditionally choose a uniform size-$S$ subset $I$. Write $J=r-S$, $h=\mathbb E J$, $\mu=r-h$, $\phi=\pi/d$, $T=nr/d$ and $D=\lfloor T\rfloor$. Define

\begin{equation*}
\begin{gathered}
g(\tau)=\mathbb E e^{-i\tau J},\quad
H_\tau=H(te^{i\tau})=H(t)e^{ir\tau}g(\tau),\\
a_s=\frac{s(d-s)}{d(d-1)},\quad
a=\mathbb E a_S,\quad a_\tau=\frac{\mathbb E[a_S e^{-i\tau J}]}{g(\tau)}.
\end{gathered}
\end{equation*}

There is an absolute $h_0>0$ such that the following holds eventually, uniformly for $0<h\le h_0$. Set

\begin{equation*}
\Psi=n\bigl(r\phi+\mathrm{Im}\log g(\phi)\bigr)
 -\frac{d-1}{2}\mathrm{Im}\log a_\phi .
\end{equation*}

At every node $\Psi=k\pi$ with integer $1\le k\le D-1$,

\begin{equation*}
(-1)^kP_n^{(r,d)}(-t^d)>0.
\end{equation*}

The logarithm of $g$ is the branch near one and that of $a_\tau$ is the branch in the right half-plane. The proof uses angular integrals throughout the near-full range, including $1\le D-k=o(d)$. It does not require asymptotic formulas for the upper coefficients.

\subsection{Tilt parameters and diagonal phases}

The successive ratios of the deficit weights decrease. The first is $q=\Pr(J=1)/\Pr(J=0)\le h/(1-h)$. Hence $\Pr(J=j)\le q^j$, $\mathbb E J^\ell=O_\ell(h)$ for every fixed $\ell\ge1$, and $h\le Cq$. Consequently $\Pr(J=0)\ge1-h$, $\Pr(J=1)\asymp h$, and $\mathrm{Var}(J)\asymp h$. Also

\begin{equation*}
a_{r-j}/a_r\le1+j,\qquad a\asymp a_r\ge 1/d.
\end{equation*}

In particular $\mathbb E[a_S^q J^\ell]=O_{q,\ell}(a^q h)$ for fixed positive integers $q,\ell$, and the same weighted estimate holds with $\mathbf1_{\{J>0\}}$ in place of $J^\ell$.

Uniformly on $[-\pi,\pi]$, $g=1+O(h)$, so $g$ stays bounded away from zero. Its fixed positive order derivatives are $O(h)$. Taylor expansion of its logarithm at zero gives $\mathrm{Im}\log g(\tau)=-h\tau+O(h|\tau|^3)$ for small $\tau$. Its centered characteristic function has second log-modulus derivative $-\mathrm{Var}(J)+O(h|\tau|)$ near zero. For far phases, the pair of atoms zero and one gives $|g(\tau)|\le\exp[-ch(1-\cos\tau)]$ by the independent-copy formula. Comparing both ranges with the value at $\phi$ gives

\begin{equation}
\log|H_\tau/H_\phi|\le -ch(\tau^2-\phi^2)\qquad(\phi\le|\tau|\le\pi),
\qquad
e^{-Ch\tau^2}\le |g(\tau)|\le1\quad (|\tau|\text{ small}).
\tag{S1}
\end{equation}

The bounds on the weighted moments of $J$ show that each of the first two derivatives of $a_\tau$ is $O(ah)$. Also $a_\tau-a=O(ah)$ uniformly in $\tau$, and its first derivative at zero is purely imaginary. Thus

\begin{equation}
\frac{\mathrm{Re}\,a_\tau}{a}=1+O(h\min(\tau^2,1)),\qquad
\frac{|\mathrm{Im}\,a_\tau|}{a}=O(h\min(|\tau|,1)),\qquad
\frac{|a_\tau|}{a}=1+O(h\min(\tau^2,1)).
\tag{S2}
\end{equation}

The cap range is chosen small enough for $\mathrm{Re}\,a_\tau$ to be positive uniformly.

\subsection{Global common-angle localization}

Use normalized Haar measure on $|z_i|=1$, $\prod_i z_i=-1$. With $R=H(t)\mathbb E\prod_{i\in I}z_i$, its integral after raising to the $n$th power equals $P_n^{(r,d)}(-t^d)$ by character orthogonality. Put $\overline z=d^{-1}\sum_i z_i$ and $D_z=\sum_i|z_i-\overline z|^2$.

At size $s$ put $m=\min(s,d-s)$. Generate a uniform subset by making $m$ disjoint random pairs from a uniform permutation, independently choosing one element from each pair, and including the next $s-m$ elements. Given the permutation, the absolute response is bounded by $e^{-X/8}$, where $X$ is the sum of squared chords of the pairs. Each chord is at most four and $\mathbb EX=2mD_z/(d-1)$. For two disjoint pairs the expected product of their squared chords is at most $d(d-1)/[(d-2)(d-3)]$ times the squared mean chord of one pair, by allowing intersections in the numerator. Thus $\mathbb EX^2\le C(\mathbb EX+(\mathbb EX)^2)$. Cauchy--Schwarz gives $\mathbb E[X/(1+X)]\ge(\mathbb EX)^2/\mathbb E[X+X^2]$ when the latter denominator is positive. With $1-e^{-x/8}\ge c x/(1+x)$ this bounds the fixed-size modulus by $1-c\min(mD_z/d,1)$, also valid when $m=0$. On averaging, retain just the size $r$ in the modulus loss. Its mass is at least $1-h$, and $\min(r,d-r)/d\ge (d-1)a_r/d$ with $a_r\asymp a$. Therefore

\begin{equation*}
|R/H(t)|\le 1-c\min(aD_z,1).
\end{equation*}

Choose a sufficiently small fixed $\eta>0$ for the local bounds below, and restrict to $D_z\le c_1\eta^2$ with a sufficiently small fixed $c_1$. The scale to compare against is $|H_\phi|^n I_0$, with

\begin{equation*}
I_0=\frac{(2\pi n a)^{-(d-1)/2}}{\sqrt d}.
\end{equation*}

On the outside part with $aD_z>1$, the integral is at most $H(t)^n e^{-cn}$. This is negligible on the principal scale since $n\gg d\log n$ under $n/d^2\to\infty$, and the principal normalization loss relative to $H(t)^n$ is $O(nh/d^2)$ in the exponent.

On the outside part with $aD_z\le1$, keep half the chord penalty for $e^{-cna\eta^2}$. Choose a pin $j$ minimizing $\sum_i|z_i-z_j|^2\le2D_z$ by averaging, and apply the other half of the penalty to the pinned squared chords. For a fixed pin, the map to the $d-1$ relative coordinates $w_i=z_i/z_j$, $i\ne j$, is onto, since one can choose $z_j^d=-1/\prod_{i\ne j}w_i$. This map is a translate of an onto homomorphism of compact tori, so it pushes normalized Haar measure to product Haar measure. Each chord factor integrates to at most $C/\sqrt{na}$ in one coordinate. A union bound over pins thus gives outside contribution at most $H(t)^n d e^{-cna\eta^2}(C/\sqrt{na})^{d-1}$. Compared to the principal scale, its logarithm is bounded by $-cna\eta^2+O(d+\log d+nh/d^2)$, tending to minus infinity because $na\gg d$ and $nh/d^2=o(na)$.

In the retained domain choose a pin minimizing $\sum_i|z_i-z_j|^2\le2D_z$ and lift all angular differences around it. Subtract their average to write $z_i=e^{i(\tau+y_i)}$, $\sum_i y_i=0$, and $\|y\|\le\eta$. The common phase $\tau$ is an odd multiple of $\phi$, taken in $(-\pi,\pi]$. This representation is unique within the small-norm domain, since any two such angle lifts differ by the same integer multiple of $2\pi$ in every coordinate, as seen by taking differences. Each chart contains a full fixed-radius ball in the zero-sum hyperplane, since $D_z\le\sum_i|e^{iy_i}-1|^2\le\|y\|^2$. In the first $d-1$ deviation coordinates the measure is $(2\pi)^{-(d-1)}dy_1\cdots dy_{d-1}$, with $y_d=-\sum_{i<d}y_i$.

\subsection{Transverse expansion and a uniform integral estimate}

Write $Q_j=\sum_i y_i^j$ and $X=\sum_{i\in I}y_i$. At a fixed size $s$, let $m_j=(s)_j/(d)_j$. Expanding the powers with $Q_1=0$ gives

\begin{equation*}
\begin{aligned}
\mathbb E[X\mid S=s]&=0,\\
\mathbb E[X^2\mid S=s]&=a_sQ_2,\\
\mathbb E[X^3\mid S=s]&=a_s\frac{d-2s}{d-2}Q_3,\\
\mathbb E[X^4\mid S=s]&=(m_1-7m_2+12m_3-6m_4)Q_4
       +3(m_2-2m_3+m_4)Q_2^2.
\end{aligned}
\end{equation*}

For example, the sums of the quartic monomials with patterns $3+1$, $2+2$, $2+1+1$ and $1+1+1+1$, with repeated exponents unlabeled, are $-Q_4$, $(Q_2^2-Q_4)/2$, $Q_4-Q_2^2/2$, and $Q_2^2/8-Q_4/4$. The coefficient of $Q_4$ is $O(a_s)$ by writing it in terms of $m_j-m_{j+1}=m_j(d-s)/(d-j)$, while $m_2-2m_3+m_4=(s)_2(d-s)_2/(d)_4=O(a_s^2)$. Consequently the real fourth moment is bounded by $C(a Q_4+a^2Q_2^2)$, using $\mathbb E a_S^2=O(a^2)$.

Tilt the size weights by $e^{-i\tau J}/g(\tau)$, whose absolute weights are bounded by a constant times the real probabilities. Since the linear term vanishes, Taylor through the cubic and the logarithm near one give

\begin{equation}
\log(R/H_\tau)
=-\tfrac12 a_\tau Q_2-\tfrac{i}{6}b_\tau Q_3
       +O(a Q_4+a^2Q_2^2),
\tag{S3}
\end{equation}

where $b_\tau$ is the tilted mean of $a_s(d-2s)/(d-2)$. The real exponential Taylor remainder is bounded by the absolute fourth moment for every real value of $X$. Also $|R/H_\tau-1|\le C a Q_2$ by the second-order Taylor bound, so the logarithm is well defined on this chart and its quadratic error is included in (S3). The weighted moment bounds give $|b_\tau|=O(a)$ and $|\mathrm{Im}\,b_\tau|=O(ah\min(|\tau|,1))$ by differentiating and by comparison with the $J=0$ term. Using $\min(|\tau|,1)\sum_i|y_i|^3\le[\min(\tau^2,1)Q_2+Q_4]/2$, the real part in (S3) is at most

\begin{equation}
-\tfrac12 a_{\rm low}Q_2+C(aQ_4+a^2Q_2^2),\qquad
a_{\rm low}=a[1-C_0h\min(\tau^2,1)]\asymp a.
\tag{S4}
\end{equation}

Fix $C_0$ large enough, then $h_0$ small enough. The same upper bound applies to the real part of the comparison exponent $-a_\tau Q_2/2$.

Here is a uniform bound for the positive exponential corrections. Put $A=na_{\rm low}\asymp na$. Under the normalized Gaussian on the hyperplane with density proportional to $e^{-A Q_2/2}$, the vector of deviations has the law of $Y_i=X_i-\overline X$ for independent $X_i\sim N(0,A^{-1})$. The vector $Y$ is independent of $\overline X$. On $Q_2(Y)\le\eta^2$, impose also $|\overline X|\le\eta$, which has probability at least $1/2$ independently. Then $|X_i|\le2\eta$, and $Q_4(Y)\le16\sum_iX_i^4$ by the fourth-power inequality and Jensen. For any fixed $C_1$, with $\eta$ sufficiently small, the one-dimensional Gaussian expectation of $e^{C_1 A X_i^4}$ restricted to this range is $1+O(A^{-1})$, by comparison with a Gaussian with half the original precision. The product is bounded since $d/A=o(1)$. This bounds the exponential moment of $C_1 n a Q_4$ on the localized ball, for any fixed constant needed below.

For the radial correction, $U=AQ_2$ has a chi-square law with $d-1$ degrees of freedom. The exponent $C_1na^2Q_2^2$ is $O(U^2/n)$. On $U\le Ld$ with fixed large $L$ this is $O(d^2/n)=o(1)$. On $Ld<U\le A\eta^2$ it is at most $\epsilon U$ for a fixed $\epsilon<1/8$ by choosing $\eta$ small. The exponential moment at $1/4$ bounds the contribution of that range by $e^{-c d}$ for $L$ sufficiently large. Cauchy--Schwarz combines the two quartic exponential bounds, and the same proof works with doubled constants.

The Gaussian for $a_{\rm low}$ has integral $I_0\exp(O(dh\min(\tau^2,1)))$, since the free-coordinate quadratic form has matrix $I+\mathbf1\mathbf1^\top$ and determinant $d$. From (S4) and the exponential bounds, the absolute contribution of any localized cluster is at most

\begin{equation}
C I_0|H_\tau|^n \exp(Cdh\min(\tau^2,1)).
\tag{S5}
\end{equation}

The exact complex Gaussian gives its leading term with an absolute error at most

\begin{equation}
C I_0|H_\tau|^n \exp(Cdh\min(\tau^2,1))\,
\left(\sqrt{\frac d{na}}+\frac d{na}+\frac{d^2}{n}\right).
\tag{S6}
\end{equation}

To see (6), apply $|e^A-e^B|\le|A-B|e^{\max(\mathrm{Re}\,A,\mathrm{Re}\,B)}$ to the two exponents and Cauchy--Schwarz to the exponential correction. Under the real comparison Gaussian, the $L^2$ norms of $Q_3,Q_4,Q_2^2$ are $O(\sqrt d\,A^{-3/2})$, $O(d A^{-2})$, and $O(d^2 A^{-2})$. The first follows from Wick's formula $\mathbb E Q_3^2=6A^{-3}(d-3+2/d)$, since the covariance has row sums zero. The others follow from marginal eighth moments and the chi-square fourth moment. Multiplying by the coefficients in (S3) gives (6). The Gaussian can be extended to the full hyperplane with an additional exponentially small error, since the chart contains a fixed-radius ball and $na\gg d$. Its integral after multiplying by the saddle value is

\begin{equation}
\frac{H_\tau^n(2\pi n a_\tau)^{-(d-1)/2}}{\sqrt d}.
\tag{S7}
\end{equation}

The factor in parentheses in (6) is $o(1)$ since $a\ge c/d$ and $n/d^2\to\infty$.

\subsection{Signs at all internal nodes}

At a node $\Psi=k\pi$, the tilted argument expansion and (S2) give

\begin{equation}
T-k=\frac{nh}{d}\left[1+O(d^{-2}+d/n)\right].
\tag{S8}
\end{equation}

Indeed $\mathrm{Im}\log g(\phi)=-h\phi+O(h\phi^3)$ and $\arg a_\phi=O(h\phi)$. The phase at $\phi$ uses $r\phi+\mathrm{Im}\log g(\phi)$ continuously, which is also the argument of $H_\phi$ between zero and $r\phi$. Since $k\le D-1$, (8) implies $nh\ge(1-o(1))d$. It also makes $T-k=O(d)$ equivalent to bounded $nh/d^2$.

If $nh/d^2\to\infty$, (S1) and (S5) make all nonprincipal clusters negligible on the principal scale. Indeed for $|\tau|\ge3\phi$ the factor $\exp(Cdh\tau^2)$ is absorbed in half the diagonal gap, because $n\gg d$ and $\tau^2\le(9/8)(\tau^2-\phi^2)$. Summing the remaining bound over odd phase labels gives $O(e^{-cnh/d^2})$. The principal conjugate pair in (7) then gives the asserted sign by (6), since its two terms have phase $\pm k\pi$.

It remains to consider a sequence on which $nh/d^2$ stays bounded. Put $L=d/\sqrt{nh}$. For a fixed large $M$ retain all odd phases $|\tau|\le M/\sqrt{nh}$. Take $M$ large enough that these include the principal phases and the discarded phases satisfy $|\tau|\ge\sqrt2\,\phi$. At a retained phase $\tau=\ell\phi$ with odd integer $\ell$, the argument of (7) can be taken as

\begin{equation*}
n\mu\tau+O(nh|\tau|^3)+O(dh|\tau|).
\end{equation*}

Subtracting $\ell$ times the principal argument gives a difference of the same error order, since $|\ell|\phi^3\le|\tau|^3$. Uniformly over retained labels it is $O_M((nh)^{-1/2}+d\sqrt{h/n})=o(1)$ by (8) and $n/d^2\to\infty$. Their signed real parts thus have the sign $(-1)^k$, because $\ell k$ has the same parity as $k$. The local errors (6) are $o(1)$ relative to the sum of their absolute leading amplitudes. In fact $dh\tau^2\le M^2d/n=o(1)$ by the retention condition, and (S2) compares the Gaussian factors uniformly to $I_0$.

On the scale $|H_\phi|^n I_0$, that positive amplitude sum is at least $c(1+L)$. If $L$ is large, use the at least $cL$ odd phases $|\tau|\le1/\sqrt{nh}$ and the lower bound in (S1). If $L$ is bounded, one principal phase suffices. For the discarded phases, (S1) and (S5) bound their total on the same scale by

\begin{equation*}
C(1+L)e^{-c'M^2}.
\end{equation*}

This follows from $\tau^2-\phi^2\ge\tau^2/2$, absorption of $dh\tau^2$ by the diagonal gap, and the Gaussian tail sum over an equally spaced lattice of spacing $2\pi\sqrt{nh}/d$ in the variable $\sqrt{nh}\,\tau$. Choosing $M$ fixed and large makes the discarded contribution smaller than the signed retained contribution. The exponential global error is already $o(1)$ on the principal scale. This proves the sign in the bounded case. Any hypothetical sequence of failed signs has a subsequence in one of these two regimes, proving the uniform eventual assertion.

\subsection{Ordering within the near-full interval}

The phase is strictly increasing with the real log-tilt $\eta_t=\log t$ in this interval for all sufficiently large parameters. Indeed

\begin{equation*}
\frac{d}{d\eta_t}\arg H_\phi
=\mathrm{Im}\frac{\mathbb E[S e^{i\phi S}]}{\mathbb E e^{i\phi S}}
=\phi\,\mathrm{Var}(J)+O(h\phi^3).
\end{equation*}

For the last equality center $J$, use its second and fourth moments in the numerator Taylor expansion, and the centered characteristic function $1+O(h\phi^2)$ in the denominator. Also $\partial_{\eta_t}a_\tau=-\mathrm{Cov}_\tau(a_S,J)$, where the covariance is under the complex size tilt. The weighted deficit moment estimates give $\partial_\tau\partial_{\eta_t}\log a_\tau=O(h)$, by differentiating this covariance and using $|a_\tau|\asymp a$. At $\tau=0$ the log-tilt derivative is real, so $(d/d\eta_t)\arg a_\phi=O(h\phi)$. Since $\mathrm{Var}(J)\asymp h$ and $n\gg d$, the derivative of $\Psi$ is positive. At $t\to\infty$ its limit is $T\pi$. Thus the near-full nodes are ordered and each target phase strictly between the initial phase on this interval and $T\pi$ is attained once.

\end{document}
